\clearpage\begingroup\SourceLeaf{601b}
\ArithmeticHead{46}{L}{PHILIPP MAENNCHEN, GAUSS ALS ZAHLENRECHNER.}
\renewcommand{\thefootnote}{\arabic{footnote})}\setcounter{footnote}{0}
\begin{center}\large IX.\par\textls[120]{Wie \textsc{Gauss} die Zahlen individualisierte.}\end{center}
\noindent\hspace*{\parindent}Was bei den zahlentheoretischen Untersuchungen fürs Zahlenrechnen
abfiel, ist für uns hier von besonderem Interesse. \textsc{Gauss} hat gelegentlich\footnote{Brief an \textsc{Schumacher} vom 6. Januar 1842, Werke XII, S. 37.}
geäussert, dass viele Zahlenrelationen ihm durch seine Beschäftigung mit der
Zahlentheorie so geläufig wären, dass sie ihm stets zur Verfügung stünden. Bei
dieser Gelegenheit führt er zwei Beispiele an: $13\mathord{.}29=377$ und $19\mathord{.}53=1007$,
\clearpage\SourceLeaf{602}\ArithmeticHead{47}{R}{IX. WIE GAUSS DIE ZAHLEN INDIVIDUALISIERTE.}
\noindent sowie die Relationen, die sich aus diesen wieder ableiten lassen. Das ist
wohl so zu verstehen, dass diese Ergebnisse durch die besondere Art, auf die
er zu ihnen gelangte, oder durch die Häufigkeit, mit der sie ihm entgegentraten, oder endlich durch den praktischen Nutzen, den sie ihm verschafften,
sich seinem Gedächtnis dauernd einprägten. Bei dem ersten Beispiel $13\mathord{.}29=377$
hat es sich vermutlich so verhalten: Diese Beziehung trat ihm wohl zum ersten
Male entgegen, als er Studien zur Berechnung der gemeinen Logarithmen
machte. Wir haben auf S. 7 schon von diesen Studien gesprochen und
brauchen daher jetzt nur auf den Bruch
\[
\frac{377}{376}=\frac{13\mathord{.}29}{2\mathord{.}2\mathord{.}2\mathord{.}47}
\]
\noindent aufmerksam zu machen. In den \textit{Tafeln zur Cyklotechnie} trat ihm dieselbe
Relation noch viel drastischer entgegen:
\[
13=3^2+4;\quad 29=5^2+4;
\]
\noindent also
\[
13\mathord{.}29=(3\mathord{.}5+4)^2+16=361+16=377.
\]
\noindent Vielleicht war es das erste Beispiel oder wenigstens eines der ersten, die er
induktiv fand, so dass es ihm besonders leicht vor Augen trat, vielleicht hat
er es auch bei der Abfassung des erwähnten Briefes aufs Geratewohl aus
seinem unerschöpflichen Vorrat herausgegriffen. Zur praktischen Verwendung
erscheint diese Relation weniger geeignet, immerhin kann man Aufgaben, wie
$77\mathord{.}377$ mit Vorteil lösen, wenn man sich der Faktorenzerlegung von 377
bedient. Denn $77\mathord{.}377=7\mathord{.}11\mathord{.}13\mathord{.}29=1001\mathord{.}29=29029$. Also kann man
man auch leicht $377^2$ bilden. Denn $377^2=300\mathord{.}377+77\mathord{.}377=111000$
$+2100+29029=142129$.

Das zweite Beispiel, $19\mathord{.}53=1007$, hat \textsc{Gauss} zweifellos wegen seines
praktischen Nutzens im Gedächtnis behalten, den es bei vielen Gelegenheiten
bringt. Denn nun kann man sofort Aufgaben lösen, wie $57\mathord{.}212=3\mathord{.}4\mathord{.}1007$
$=12084$ und andere. Auch diese Relation ist ihm vermutlich zuerst bei den
gemeinen Logarithmen entgegengetreten, nämlich bei der Betrachtung des
Bruches $\frac{1008}{1007}=\frac{2^4\mathord{.}3^2\mathord{.}7}{19\mathord{.}53}$. Ein anderer Ausgangspunkt für diese Relation mag
\clearpage\SourceLeaf{603}\ArithmeticHead{48}{L}{PHILIPP MAENNCHEN, GAUSS ALS ZAHLENRECHNER.}
\setcounter{footnote}{0}
\noindent die folgende Überlegung gewesen sein\footnote{Wahrscheinlich ist ihm eine ganze Gruppe von Multiplikationsresultaten bekannt gewesen, die in
der Nähe von 1000 liegen. $42\mathord{.}48=2016$; also $21\mathord{.}48=42\mathord{.}24=12\mathord{.}84=36\mathord{.}28=1008$; $44\mathord{.}46$
$=2024$; also $22\mathord{.}46=44\mathord{.}23=1012$. Ferner $51\mathord{.}59=3009$; $17\mathord{.}59=1003$. Die übrigen zusammengesetzten Zahlen in der Nähe von 1000 müssen eben planmässig zerlegt werden, falls nicht wieder besondere Kunstgriffe möglich sind. — Auch in der Nähe der Million hat \textsc{Gauss} solche Zerlegungen gesucht,
wie eine auf dem Deckblatt von \textsc{Lamberts} Tafeln befindliche Aufzeichnung: \textit{Analysis numerorum supra
millionem} zeigt.}:
\[
57\mathord{.}53=50\mathord{.}60+21=3021,
\]
\noindent also
\[
19\mathord{.}53=1007.
\]
\noindent Wie wir auf S. 27 sahen, war ihm auch $29\mathord{.}41=1189$ geläufig. Auch diese
Relation trat ihm bei verschiedenen Gelegenheiten entgegen: $(30-1)(40+1)$
$=1200-20+9$ oder $29\mathord{.}41=35^2-6^2$, oder endlich in der Cyklotechnie als
Sonderfall der Formel
\[
(a^2+b^2)(c^2+d^2)=(ac+bd)^2+(ad-bc)^2.
\]
\noindent Danach ergibt sich
\[
29\mathord{.}41=(5^2+2^2)(5^2+4^2)=33^2+10^2=30^2+17^2=1189.
\]
\noindent Dieses Individualisieren, in Verbindung mit der welschen Praktik, kam ihm
auch bei vielen Aufgaben des schriftlichen Rechnens, namentlich bei Multiplikationsaufgaben, zustatten. Hätte er etwa die Aufgabe zu lösen gehabt
(vergl. \textsc{Remer}, a. a. O., S. 120 ff.):
\[
423219\mathord{.}272673,
\]
\noindent so würde er, sofort erkennend, dass $2673=2700-27=27\mathord{.}99$ ist, so gerechnet haben:
\[
\begin{array}{r@{\quad}l}
423219\mathord{.}27\hphantom{0}&(27=9\mathord{.}3)\\\cline{1-1}
3808971\hphantom{0000}&\\\cline{1-1}
114269130000&[x\mathord{.}270000]\\
1142691300&[x\mathord{.}2700]\\\cline{1-1}
115411821300&[x\mathord{.}272700]\\
11426913&[-x\mathord{.}27]\\\cline{1-1}
115400394387.&
\end{array}
\]
\clearpage\SourceLeaf{604}\ArithmeticHead{49}{R}{X. CHRONOLOGISCHE ARBEITEN.}
\setcounter{footnote}{0}
\noindent Man vergleiche damit die gewöhnliche Ausführung:
\[
\begin{array}{l}
\hphantom{0}423219\mathord{.}272673\\\hline
\hphantom{0}846438\\
\hphantom{0}2962533\\
\hphantom{000}846438\\
\hphantom{000}2539314\\
\hphantom{0000}2962533\\
\hphantom{00000}1269657\\\hline
115400394387.
\end{array}
\]
\noindent Dieses Beispiel wird uns später noch einmal beschäftigen.

Wie schon wiederholt bemerkt wurde, haben sich viele der von \textsc{Gauss}
untersuchten Zahlen (Summen von Reihen, Logarithmen, Wurzeln, reziproke
Werte, arithmetisch-geometrische Mittel) infolge der allseitigen Behandlung
vieler Zahlen, Zahlengruppen und Zahlenrelationen seinem schon von Natur
ungewöhnlichem Gedächtnis dauernd eingeprägt. Und wenn er bei der Auswertung einer Konstanten auf eine schon bekannte Zahl stiess, so bemerkte
er das wahrscheinlich schon, ohne eine Tabelle nachzusehen. Das berühmteste
Beispiel ist wohl das arithmetisch-geometrische Mittel von 1 und $\sqrt{2}$, dessen
Wert sich\footnote{\textit{Tagebuch} Nr. 98, Werke X, 1, S. 542.} als gleich erwies dem Wert $\frac{\pi}{\tilde{\omega}}$, wo $\tilde{\omega}$ die Konstante ist, die bei
den lemniskatischen Funktionen die entsprechende Rolle spielt, wie $\pi$ bei den
Kreisfunktionen.
\clearpage\endgroup
